\section{Methodology}

CricketMind is built as one connected pipeline that can answer two different kinds of cricket questions: a question about a specific video clip of a shot, and a question that is only text, with no video at all. When a user submits a query, the system first checks whether a video clip is attached. If a clip is attached, the video is uniformly sampled into fifteen frames, since this is enough to capture the whole shot, from the backlift to the follow-through, without wasting effort on repeated near-identical frames, and each frame is passed through a vision encoder that turns it into a numeric description of the batter's posture, bat position, and ball location. At the same time, and on both paths, the text of the query is turned into a numeric representation by a separate text encoder, so the system always understands which player, shot, or situation is being asked about, even when there is no video. When a clip is present, the fifteen frame descriptions are combined by a small temporal transformer that lets the very first frame of the shot directly influence the very last frame, which is what allows the system to tell apart shots that look similar frame by frame but differ mainly in the order and speed of motion, such as a sweep and a reverse sweep; the combined description is then passed through a classifier that predicts one of fifteen shot types along with a confidence score, and if this confidence is too low the system does not accept the result but instead tries again with a different camera angle or a different set of frames, up to a limited number of attempts, so an unreliable label is never passed forward silently. The accepted shot type, together with the query and extra posture measurements such as weight transfer and bat angle, is then combined into one structured record called the Shot DNA; a text-only query skips the entire video path and arrives at this same structured record directly, carrying only the information taken from its own text, and in that case the posture measurements are simply left out rather than filled in with a guess, because inventing a measurement that was never actually taken would quietly poison the very record that the rest of the system is built to trust. This structured record is then used to search two separate databases: a precise ball-by-ball database that returns very specific matching evidence, and a broader historical database that returns general career trends, so the final analysis is neither too narrow nor too generalized, and two shared calculation engines take over from there, one that computes exact numbers such as strike rate and dismissal rate from the retrieved match data, and another that answers "what if" questions about future situations by running many randomized simulated outcomes and reporting how likely each result is; these two engines are the only parts of the system allowed to produce a number, so every other component has to borrow numbers from them rather than making up its own. A routing step then looks at the structured record, the question, and the computed numbers, and wakes up only the small group of specialized agents that are relevant to that question, out of twelve total agents built for different cricket stakeholders such as coaches, players, scouts, medical staff, and broadcasters, and on a text-only query the router also skips any agent, such as the ones focused on personal technique or player development, whose answer would depend on posture measurements that were never actually taken, which keeps the final report shorter but still fully backed by real evidence. Before any agent's answer is shown to the user, every number it mentions is checked against the numbers that were actually produced by the two calculation engines, and if an agent states a number that was never really computed, the system asks that agent to redo its answer using the real evidence instead, so no invented number can reach the user; finally, all the checked and approved agent answers are combined into one single, readable report that directly answers the user's original question.